% Required in the paper preamble for the boxed user input example:
% \usepackage[most]{tcolorbox}

\subsection{Starting a new session}

After preparing the input statement and data files, the session directory should
contain the user-supplied files:

\begin{verbatim}
sessions/<session-name>/
  inputs/
    user_info.txt
    <data-file>.csv
    ...
\end{verbatim}

The file \texttt{user\_info.txt} is the reviewed scientific description of the
problem, and the CSV files are the experimental records referenced by that
description. The generated YAML and Python files are not starting inputs. They
are created by the workflow.

Create the session with:

\begin{verbatim}
pfit new sessions/<session-name>
\end{verbatim}

The configured local model extracts the problem in bounded stages using the
input statement and CSV headers. The stages identify the selected experiments,
trainable and fixed parameters, integrated states and initial conditions,
equations, derived observables, measured and forcing columns, and the loss
definition. If a required item is missing or contradictory, the command reports
targeted \texttt{missing\_inputs} and stops. It does not conduct an interactive
question-and-answer session or silently invent missing equations, parameter
bounds, initial conditions, data mappings, units, or loss terms.

\paragraph{Running example.}
The tutorial uses a two-state drug-response model fitted simultaneously to pulse
and washout experiments. Both CSV files use the same columns,

\begin{verbatim}
time,dose_rate,signal
\end{verbatim}

but have different forcing histories and initial conditions. The corresponding
\texttt{user\_info.txt} is:

\begin{tcblisting}{
  listing only,
  breakable,
  colback=gray!6,
  colframe=gray!35,
  boxrule=0.4pt,
  arc=1mm,
  left=2mm,
  right=2mm,
  top=1mm,
  bottom=1mm
}
Problem:
Fit pulse.csv and washout.csv using one shared parameter set.

Experiments and initial conditions:
pulse.csv:   D(0) = 0, R(0) = 0
washout.csv: D(0) = 0, R(0) = 0.3

CSV columns:
time: time in minutes
dose_rate: known forcing input in concentration per minute;
           use linear interpolation
signal: measured value of signal_model in arbitrary units

Trainable parameters:
k_elim:   bounds [0.01, 10], logarithmic
k_on:     bounds [0.001, 10], logarithmic
k_off:    bounds [0.001, 10], logarithmic
baseline: bounds [50, 150], linear
gain:     bounds [1, 200], linear

Fixed parameters:
None.

ODE system:
dD/dt = dose_rate - k_elim*D
dR/dt = k_on*D*(1 - R) - k_off*R

Observable:
signal_model = baseline + gain*R
CSV column signal measures signal_model.

Loss:
For each experiment e, use only finite signal measurements.

scale_e = max(abs(finite signal values in experiment e))

L_e = sqrt(mean(
    ((signal_model_e - signal_e)/scale_e)^2
))

The total loss is the equal-weight mean:
L = (L_pulse + L_washout)/2.

Outputs:
Write time, D, R, measured signal, and signal_model.
\end{tcblisting}

This statement supplies the scientific decisions required by the workflow. The
local model is not expected to choose missing equations, parameter bounds,
initial conditions, data mappings, units, or loss terms. The measured CSV column
\texttt{signal} is valid here because the statement explicitly maps it to the
derived observable \texttt{signal\_model}. In general, measured columns should
map unambiguously to declared integrated states or derived observables, and
forcing or uncertainty columns should be identified by role.

\paragraph{Generated files.}
If extraction and validation succeed, \texttt{pfit new} creates:

\begin{verbatim}
sessions/<session-name>/
  inputs/
    user_info.txt
    user_input.yaml
    pulse.csv
    washout.csv
  generated/
    user_model.py
    agent_logs/
\end{verbatim}

The file \texttt{inputs/user\_input.yaml} contains the structured experiment
definitions, column roles, initial conditions, parameter bounds, state ordering,
numerical settings, and output configuration. Fresh generated configurations
also include provisional numerical automation settings used later by
\texttt{pfit jax}, including automatic integrator selection, automatic
\texttt{max\_steps} estimation, bounded solver-step recovery, tolerance
calibration, and solver-accuracy checks. These settings are subject to later
deterministic validation and may be adjusted during JAX generation.

The file \texttt{generated/user\_model.py} contains the readable NumPy
implementation of the ODEs, forcing interpolation, observables, loss, and output
calculation. Detailed model calls and workflow events are stored under
\texttt{generated/agent\_logs/}; an additional
\texttt{pfit\_new\_review.txt} is written when extraction produces review notes.

Before these files are accepted, the framework checks the CSV headers, declared
names, numeric values, parameter bounds, initial conditions, equation symbols,
column mappings, generated Python interfaces, and selected source-consistency
conditions. Fixed values and explicitly stated initial conditions are checked
against the original problem statement. Bounded repair attempts may correct
eligible expression or response-format errors, but deterministic validation
still decides acceptance.

Passing \texttt{pfit new} means that the generated artifacts satisfy the checks
implemented for new-session extraction. It does not prove that the extracted
model represents the intended science, that the future JAX translation is
correct, or that the solver settings are adequate for fitting. Before
continuing, the user should compare the generated files with the original
specification and verify the selected datasets, parameter bounds, state
ordering, initial conditions, equations, forcing interpolation, observable
mappings, loss, output definitions, and numerical settings.

After intentionally changing the source statement, regenerate the session with:

\begin{verbatim}
pfit new sessions/<session-name> --overwrite
\end{verbatim}

This replaces the generated new-session artifacts while preserving the
user-supplied \texttt{user\_info.txt} and CSV files.
